\begin{longtable}{
    %|>{\raggedright\arraybackslash}p{0.27\textwidth}
    |>{\raggedright\arraybackslash}p{0.23\textwidth}
    |p{0.55\textwidth}
    |p{0.32\textwidth}
    |p{0.27\textwidth}|
    } % 5 - 1.35 \ 
\caption{Классификация терминов цикла по уровню анализа экономики}
\label{tab:cycle-classification-level} \\
\LongTableHeader{
    \hline
    \rowcolor{black}
    %\textcolor{white}{\textbf{Русский термин}} &
    \textcolor{white}{\textbf{Термин}} &
    \textcolor{white}{\textbf{Описание}} &
    \textcolor{white}{\textbf{Источник}} &
    \textcolor{white}{\textbf{Цитата из источника}} \\
    \hline
}

\longcaptioneachpage{Источник: составлено автором}

% Национальный цикл
\begin{minipage}[t]{\linewidth}
    \raggedright
    \begin{itemize}
        \item \textbf{национальный цикл (рус.)}
        \item national cycle (англ.)
        \item cycle national (фр.)
        \item nationaler konjunkturzyklus (нем.)
    \end{itemize}
\end{minipage}
&
Базовый и наиболее изученный уровень в иерархии циклического анализа: единицей наблюдения служит экономика суверенного государства в целом. Burns \& Mitchell (1946) определили эталонный \textquotedbl{}reference cycle\textquotedbl{} именно на национальном уровне, опираясь на агрегированные показатели деловой активности США; NBER официально датирует пики и впадины американского национального цикла по сей день. Gallarotti (2001) дал краткую энциклопедическую характеристику: национальный цикл – \textquotedbl{}систематический путь\textquotedbl{}, которым совокупная экономическая активность нации следует в течение определённого периода. Ключевая внутренняя проблема уровня: различные регионы и отрасли внутри одной страны могут находиться в разных фазах цикла, что ставит вопрос о репрезентативности агрегированного индикатора. В современных исследованиях синхронизации (Lopes et al., 2018) национальные циклы 12 и более стран сопоставляются на временных горизонтах в несколько десятилетий, что позволяет выявить устойчивые кластеры и расхождения. Начиная с 1990-х гг. анализ национального цикла обогащается факторными моделями, позволяющими отделить \textquotedbl{}мировую, региональную и страновую\textquotedbl{} компоненты.
&
\begin{minipage}[t]{\linewidth}
    \vspace{-7pt}
    \RaggedRight
    \begin{enumerate}
        \item Burns A. F., Mitchell W. C. (1946). \textit{Measuring Business Cycles}\cite{Burns-1946}
        \item Gallarotti G. M. (2001). Business Cycles ?
    \end{enumerate}
\end{minipage}
&
\textquotedbl{}\textquotedbl{}
\begin{flushright}
    
\end{flushright}
\\ \hline

\end{longtable}